Almost half of the stars that we see are either binary or multiple star
systems. A well-known example of this is Menkalinan (Beta Aurigae), which was
initially thought to be a single star, but today recognised as a binary
system comprising two stars that we will refer to as Menkalinan A and B. In
the following figure, a spectrum of the system (obtained by the observatory
of the Universidad de los Andes, in Bogot\'a) is shown: a spectrum of the
Menkalinan binary system in the region of $H_\alpha$; the $y$-axis is for the
relative flux, and the $x$-axis measures wavelengths. Menkalinan A is marked
as A in the graph, and Menkalinan B is marked as B.

Answer the following questions using the plot and noting that the wavelength
of the $H_\alpha$ line in the laboratory frame is 656.28\,nm. Assume circular
orbits, and assume that the binary system as a whole is at rest with respect
to the observer.

\begin{description}
\item[(7.1)] In the spectrum, we can see the $H_\alpha$ line for each star
  in the system. Calculate the line-of-sight velocity of each star (km/s)
  and determine, at the time of this observation, which of the two stars is
  moving towards us. \hfill[5.0\,pt]
\item[(7.2)] The binary system is located 81.1 light years from Earth and
  has an orbital period of 3.96 days. The semi-major axis for Menkalinan B
  (smaller star) was measured to be 3.35 milliarcseconds. If the mass ratio
  of the two components is 1.026, find the total mass of the system (in
  solar masses). \hfill[4.0\,pt]
\item[(7.3)] Calculate the individual masses of Menkalinan A and B in solar
  masses. \hfill[2.0\,pt]
\item[(7.4)] Since Menkalinan A and B are main sequence stars, use the
  relation:
  \[ \frac{L}{L_\odot} = \left(\frac{M}{M_\odot}\right)^{3.5} \]
  to estimate the luminosity of each star (in solar luminosity).
  \hfill[2.0\,pt]
\end{description}
